\documentclass[aspectratio=169, 12pt]{beamer}

\usetheme{Madrid}
\usecolortheme{whale}

\usepackage[utf8]{inputenc}
\usepackage[T1]{fontenc}
\usepackage[italian]{babel}
\usepackage{amsmath,amssymb,amsfonts}
\usepackage{mathtools}
\usepackage{booktabs}

\newcommand{\CL}[1]{\llbracket #1 \rrbracket}

\title{Coalition Logic and Quantification}
\subtitle{From Modal Logic to Quantified Coalition Logic}
\author{Luca Sforza}
\date{\today}
\institute{Logic Course --- Final Presentation}

\begin{document}

% ============================================================
% SLIDE 1 — Title
% ============================================================
\begin{frame}
  \titlepage
\end{frame}

\note{
  + Welcome everyone. State your name and the topic.
  + This presentation covers coalition logic and its quantified extension.
  + The talk is structured in 8 parts, lasting about 30 minutes.
  + You will use game-theoretic frameworks to reason about multi-agent systems.
  + Timing: 0:00 to 0:30.
}

% ============================================================
% SLIDE 2 — Outline
% ============================================================
\begin{frame}{Outline}
  \begin{enumerate}
    \item Introduction: From Kripke Models to Game Frames
    \item The Concept of Coalitional Power
    \item Coalition Logic (CL): Syntax and Semantics
    \item Axiomatization and Limits of CL
    \item Quantified Coalition Logic (QCL)
    \item Quantified Operators
    \item Metatheory: Computational Complexity
    \item Applications: Voting and Social Choice
    \item Conclusions and Reflections
  \end{enumerate}
\end{frame}

\note{
  + Briefly walk through the 8 sections.
  + Emphasize the logical progression: single-agent modal logic, multi-agent, coalition logic, quantified coalition logic.
  + Mention that Section 7 (complexity) and Section 8 (voting applications) tie everything together.
  + Timing: 0:30 to 1:00.
}

% ============================================================
% SLIDE 3 — Introduction: From Kripke to Game Frames
% ============================================================
\section{Introduction: From Kripke Models to Game Frames}

\begin{frame}{Single-Agent Modal Logic}
  \begin{itemize}
    \item In the simplest case, we have \textbf{one agent} who can choose between actions.
    \item A \textbf{Kripke model} uses an accessibility relation $R$ that associates to every state all reachable states.
    \item A modal formula $\Diamond \varphi$ means: ``the agent can act so that $\varphi$ will be true.''
  \end{itemize}
\end{frame}

\note{
  + Start with the familiar: Kripke models for single-agent modal logic.
  + An agent at state $s$ can move to state $t$ if $s R t$.
  + The modality $\Diamond \varphi$ expresses possibility: there exists an accessible state where $\varphi$ holds.
  + This is the foundation that we will generalize.
  + Timing: 1:00 to 2:00.
}

\begin{frame}{The Multi-Agent Problem}
  \begin{itemize}
    \item Extending to multiple agents with individual accessibility relations $R_i$ considers agents \textbf{in isolation}.
    \item But agents' actions \textbf{interact}: what happens when agent 1 and agent 2 act \textbf{simultaneously}?
    \item The outcome depends on the \textbf{combination} of all participants' choices.
  \end{itemize}
\end{frame}

\note{
  + This is the key motivation: standard multi-agent logics treat agents independently.
  + But in reality, actions interact. If agent 1 moves left and agent 2 moves right, the result is a joint outcome.
  + Game frames solve this by associating each state with a strategic game.
  + The outcome function maps strategy profiles to states.
  + This is the foundation of Coalition Logic.
  + Timing: 2:00 to 4:00.
}

\begin{frame}{The Solution: Game Frames}
  \begin{itemize}
    \item Marc Pauly introduces \textbf{game frames} (2002).
    \item Each state is associated with a \textbf{strategic game} among the agents.
    \item Outcomes of the game are states of the model, leading to \textbf{Coalition Logic} (CL).
  \end{itemize}
\end{frame}

% ============================================================
% SLIDE 5 — The Concept of Coalitional Power
% ============================================================
\section{The Concept of Coalitional Power}

\begin{frame}{Strategic Games}
  A strategic game $G = (N, \{\Sigma_i \mid i \in N\}, o, S)$ consists of:
  \begin{itemize}
    \item $N$: finite set of agents (players)
    \item $\Sigma_i$: set of strategies for agent $i$
    \item $S$: set of outcome states
    \item $o: \prod_{i \in N} \Sigma_i \rightarrow S$: outcome function
  \end{itemize}
\end{frame}

\note{
  + Define the formal building block: a strategic game.
  + $N$ is the set of players, $\Sigma_i$ the strategies, $S$ the states.
  + The outcome function $o$ takes a strategy profile (one strategy per player) and returns a state.
  + This is a game form with no preferences yet, just the mechanics.
  + Timing: 4:00 to 5:00.
}

\begin{frame}{Game Frames}
  \begin{block}{Definition}
    Let $\Gamma^N_S$ be the set of all strategic games over agents $N$ and states $S$.
    A \emph{game frame} is a pair $(S, \gamma)$ where $\gamma: S \rightarrow \Gamma^N_S$ assigns a strategic game to each state.
  \end{block}

  \begin{itemize}
    \item Game frames are essentially \textbf{extensive games with simultaneous moves}.
    \item At every state, agents play a game; the outcome leads to a new state (and a new game).
  \end{itemize}
\end{frame}

\note{
  + A game frame is a transition system where transitions are strategic games.
  + At each state $s$, the function $\gamma$ gives the game that agents must play.
  + The outcome of that game determines the next state.
  + This is more general than Kripke frames: instead of a simple relation, we have a full game.
  + Timing: 5:00 to 6:00.
}

\begin{frame}{From Strategy to Power: Effectivity Functions}
  \begin{block}{Definition}
    The \emph{effectivity function} $E_G: \mathcal{P}(N) \rightarrow \mathcal{P}(\mathcal{P}(S))$ derived from game $G$ is defined as:
    \[ X \in E_G(C) \iff \exists \sigma_C \forall \sigma_{N \setminus C}: o(\sigma_C, \sigma_{N \setminus C}) \in X \]
  \end{block}

  \begin{itemize}
    \item Coalition $C$ is \textbf{effective} for $X$ if $C$ has a joint strategy that \textbf{forces} the outcome into $X$, regardless of what the others do.
  \end{itemize}
\end{frame}

\note{
  + This is the central concept: effectivity functions.
  + $X \in E_G(C)$ means coalition $C$ can force the outcome into set $X$.
  + The quantifier structure is key: EXISTS a strategy for $C$, FOR ALL strategies of the opponents.
  + This captures the idea of ``forcing'': the coalition has a winning strategy.
  + The logic will reason about these effectivity functions directly.
  + Timing: 6:00 to 8:00.
}

\begin{frame}{Playability Properties}
  Not every function can represent a real game. An effectivity function must be \textbf{playable}:

  \begin{description}
    \item[Outcome-monotonicity] If $C$ can force $X$, it can force any $Y \supseteq X$.
    \item[Superadditivity] If disjoint $C_1, C_2$ force $X_1, X_2$ respectively, then $C_1 \cup C_2$ can force $X_1 \cap X_2$.
    \item[$N$-maximality] The grand coalition can force $X$ iff the empty coalition cannot force $S \setminus X$.
  \end{description}
\end{frame}

\note{
  + These three properties ensure the effectivity function comes from a real game.
  + Outcome-monotonicity: if you can guarantee winning, you can guarantee ``someone wins.''
  + Superadditivity: disjoint coalitions can combine their strategies, cooperation pays off.
  + $N$-maximality: if nobody can prevent $X$, then everyone together can achieve it.
  + These properties will be reflected as axioms in the logic.
  + Timing: 8:00 to 9:00.
}

% ============================================================
% SLIDE 9 — Coalition Logic: Syntax
% ============================================================
\section{Coalition Logic (CL): Syntax and Semantics}

\begin{frame}{Syntax: The Coalition Operator}
  Given agents $\mathcal{A}g$ and atomic propositions $\Phi$:
  \[ \varphi ::= p \mid \neg \varphi \mid \varphi_1 \lor \varphi_2 \mid \langle\!\langle C \rangle\!\rangle \varphi \]

  \begin{itemize}
    \item $\langle\!\langle C \rangle\!\rangle \varphi$ means: ``coalition $C$ has a joint strategy to \textbf{force} an outcome satisfying $\varphi$, regardless of the actions of $\mathcal{A}g \setminus C$.''
  \end{itemize}
\end{frame}

\note{
  + Introduce the modal operator of Coalition Logic.
  + The double angle brackets denote the coalition ability operator.
  + It is read as ``C can achieve phi.''
  + We use this notation to avoid confusion with the universal quantifier $[P]\varphi$ in QCL.
  + Timing: 9:00 to 10:00.
}

\begin{frame}{From Coalition Frames to Models}
  \begin{description}
    \item[Coalition Frame] A pair $\mathcal{F} = (S, E)$ where $E$ assigns a playable effectivity function to each state.
    \item[Coalition Model] A triple $\mathcal{M} = (S, E, \pi)$ adding a valuation $\pi$.
  \end{description}

  \vspace{0.5em}
  The semantics:
  \[ \mathcal{M}, s \models \langle\!\langle C \rangle\!\rangle \varphi \iff \varphi^{\mathcal{M}} \in E(s, C) \]
  where $\varphi^{\mathcal{M}} = \{t \in S \mid \mathcal{M}, t \models \varphi\}$.
\end{frame}

\note{
  + Frames give structure; models add truth values for propositions.
  + The semantics is elegant: $C$ can force $\varphi$ iff the set of $\varphi$-states is in the effectivity of $C$ at $s$.
  + This directly connects the logic to the game-theoretic notion of effectivity.
  + The formula $\varphi^{\mathcal{M}}$ is the set of states where $\varphi$ is true.
  + Timing: 10:00 to 11:00.
}

\begin{frame}{The Characterization Theorem (Theorem 3.2)}
  \begin{block}{Theorem (Characterization)}
    An effectivity function $E$ is playable if and only if there exists a strategic game $G$ such that the effectivity function induced by $G$ coincides with $E$.
  \end{block}
\end{frame}

\note{
  + This is the most important theoretical result in CL.
  + It bridges two worlds: the abstract algebraic world of effectivity functions, and the concrete world of strategic games.
  + The forward direction is straightforward: every game induces a playable effectivity function.
  + The reverse direction is the deep one: every playable function comes from some game.
  + This justifies working directly with effectivity functions in the logic.
  + Timing: 11:00 to 12:00.
}

\begin{frame}{Proof Sketch}
  $(\Rightarrow)$: If $E$ is induced by a game $G$:
  \begin{itemize}
    \item Outcome-monotonicity: same strategy works for supersets.
    \item Superadditivity: disjoint strategies do not interfere.
    \item $N$-maximality: the grand coalition controls everything.
  \end{itemize}

  \vspace{0.5em}
  $(\Leftarrow)$: If $E$ is playable, construct an artificial game:
  \begin{itemize}
    \item Introduce strategies representing guaranteed choices from $E$.
    \item Define outcomes so that a winning strategy produces exactly the declared set.
    \item Playability properties ensure consistency.
  \end{itemize}
\end{frame}

\note{
  + The forward direction is almost immediate from the definitions.
  + For superadditivity: if $C_1$ and $C_2$ are disjoint, their strategies do not conflict, so they can be played simultaneously.
  + For the reverse direction: the construction is non-trivial. You build a game where each coalition's ``power'' matches what $E$ declares.
  + The three playability properties are essential for the construction to be consistent.
  + Timing: 12:00 to 14:00.
}

% ============================================================
% SLIDE 13 — Axiomatization
% ============================================================
\section{Axiomatization and Limits of CL}

\begin{frame}{Axiom System}
  The axiomatization by Pauly (complete for weak playability models):

  \begin{description}
    \item[Consistency ($\bot$)] $\neg \langle\!\langle C \rangle\!\rangle \bot$ --- no coalition can force the impossible.
    \item[Efficacy ($\top$)] $\langle\!\langle C \rangle\!\rangle \top$ --- every coalition can always force something trivially true.
    \item[$\mathcal{A}g$-maximality] $\neg \langle\!\langle \emptyset \rangle\!\rangle \neg \varphi \to \langle\!\langle \mathcal{A}g \rangle\!\rangle \varphi$.
    \item[Superadditivity ($S$)] $(\langle\!\langle C_1 \rangle\!\rangle \varphi_1 \land \langle\!\langle C_2 \rangle\!\rangle \varphi_2) \to \langle\!\langle C_1 \cup C_2 \rangle\!\rangle (\varphi_1 \land \varphi_2)$ for $C_1 \cap C_2 = \emptyset$.
  \end{description}
\end{frame}

\note{
  + These axioms mirror the playability properties.
  + Consistency: nothing can force a contradiction.
  + Efficacy: every coalition can always do something.
  + $\mathcal{A}g$-maximality: what the empty coalition cannot make inevitable, the grand coalition can achieve.
  + Superadditivity is the most critical: it captures the ability to combine strategies.
  + This axiom is what makes satisfiability PSPACE-complete rather than NP-complete.
  + Timing: 14:00 to 15:00.
}

\begin{frame}{The Succinctness Problem}
  To express ``some coalition can force $\varphi$'' in CL:
  \[ \langle\!\langle C_1 \rangle\!\rangle \varphi \lor \langle\!\langle C_2 \rangle\!\rangle \varphi \lor \dots \lor \langle\!\langle C_{2^n} \rangle\!\rangle \varphi \]

  \begin{itemize}
    \item With $n$ agents, this formula grows \textbf{exponentially} ($2^n$ disjuncts).
    \item Expressing general properties like ``a majority can win'' is \textbf{practically impossible} in CL.
  \end{itemize}
\end{frame}

\note{
  + This is the key limitation that motivates QCL.
  + In CL, you must enumerate every possible coalition explicitly.
  + For $n$ agents, there are $2^n$ coalitions, exponential blowup.
  + A simple property like ``some majority can force $\varphi$'' requires listing all majority coalitions.
  + For $n = 100$ voters, this is astronomically large.
  + Timing: 15:00 to 17:00.
}

% ============================================================
% SLIDE 15 — QCL: Coalition Predicates
% ============================================================
\section{Quantified Coalition Logic (QCL): Coalition Predicates}

\begin{frame}{The Innovation: Predicates over Coalitions}
  Instead of enumerating members, describe coalitions via \textbf{predicates} $P$:
  \begin{itemize}
    \item \textbf{Inclusion:} Coalition $C$ is a subset of $D$.
    \item \textbf{Containment:} Coalition $C$ contains $D$.
    \item Combined with Boolean connectives.
  \end{itemize}

  \vspace{0.5em}
  \textbf{The Cardinality Predicate:}
  \begin{itemize}
    \item A predicate that identifies coalitions of size at least $n$.
    \item Written as $|C| \geq n$.
  \end{itemize}
\end{frame}

\note{
  + QCL's innovation is not more expressive power, but more concise representation.
  + Predicates describe properties of coalitions: is this coalition a subset of $D$? Does it contain $D$?
  + The cardinality predicate is crucial: it identifies all coalitions of size at least $n$.
  + Majority becomes a single predicate instead of an exponential disjunction.
  + Timing: 17:00 to 19:00.
}

\begin{frame}{Quantified Operators: $\langle P \rangle$ and $[P]$}
  \begin{block}{Existential}
    $\langle P \rangle \varphi$ --- there exists a coalition $C$ satisfying $P$ that can force $\varphi$.
  \end{block}

  \begin{block}{Universal}
    $[P] \varphi$ --- all coalitions satisfying $P$ can force $\varphi$.
  \end{block}

  \vspace{0.5em}
  Semantics:
  \[ \mathcal{M}, s \models [P] \varphi \iff \forall C \subseteq \mathcal{A}g: (C \models P \implies \varphi^{\mathcal{M}} \in E(C, s)) \]
\end{frame}

\note{
  + The existential $\langle P \rangle \varphi$: SOME coalition with property $P$ can force $\varphi$.
  + The universal $[P] \varphi$: ALL coalitions with property $P$ can force $\varphi$.
  + Important: $[P]$ is NOT simply the dual of $\langle P \rangle$. The dual would be $\forall \exists$, but $[P]$ follows $\forall \exists \forall$, the correct quantifier pattern for collective ability.
  + This subtlety is essential for modeling coalition power correctly.
  + Timing: 19:00 to 21:00.
}

\begin{frame}{Exponential Succinctness of QCL}
  \begin{block}{Theorem (Succinctness)}
    There exists an infinite family of QCL formulas such that any equivalent CL formula has exponential length.
  \end{block}

  \vspace{0.5em}
  Example: ``some majority can force $\varphi$''

  \begin{itemize}
    \item \textbf{In QCL:} $\langle \text{majority} \rangle \varphi$ --- one formula, constant size.
    \item \textbf{In CL:} exponential disjunction of all majority coalitions.
  \end{itemize}
\end{frame}

\note{
  + This is the main theoretical result of the paper by Agotnes, van der Hoek, and Wooldridge.
  + QCL and CL have the same expressive power: nothing new is expressible.
  + But QCL is exponentially more succinct: the same property can be stated with much shorter formulas.
  + The proof compares the length of formulas in both logics for the majority property.
  + In CL, you must list all majority coalitions. In QCL, just one predicate.
  + Timing: 21:00 to 23:00.
}

% ============================================================
% SLIDE 18 — Complexity
% ============================================================
\section{Metatheory: Computational Complexity}

\begin{frame}{CL is PSPACE-complete}
  \begin{itemize}
    \item Satisfiability for CL is \textbf{PSPACE-complete} (Pauly 2002).
    \item This is harder than propositional logic (NP-complete) because of the \textbf{superadditivity axiom}.
    \item The computer must reason about how strategies \textbf{combine} when coalitions form.
  \end{itemize}
\end{frame}

\note{
  + The superadditivity axiom is what makes CL harder than propositional logic.
  + You cannot just check each coalition independently; you must consider how they combine.
  + The surprise: despite being exponentially more succinct, QCL has the SAME complexity as CL.
  + This means QCL is strictly better for humans (shorter formulas) without being worse for machines.
  + This is a rare and desirable property in logic.
  + Timing: 23:00 to 26:00.
}

\begin{frame}{QCL: No Additional Cost}
  \begin{itemize}
    \item \textbf{Satisfiability:} PSPACE-complete, same as CL.
    \item \textbf{Model checking:} polynomial time (explicit models), PSPACE-complete (concise representations like RML).
  \end{itemize}

  \vspace{0.5em}
  Despite being exponentially more succinct, QCL has the \textbf{same} complexity as CL.
\end{frame}

% ============================================================
% SLIDE 19 — Application: Majority Voting
% ============================================================
\section{Applications: Voting and Social Choice Mechanisms}

\begin{frame}{Majority Voting: CL vs QCL}
  Consider $n$ agents choosing between outcomes $\omega_1$ and $\omega_2$.

  \vspace{0.5em}
  \textbf{In QCL, compact specification:}
  \begin{itemize}
    \item Majority power: $[\text{majority}] \omega_1 \land [\text{majority}] \omega_2$
    \item Minority restriction: $(\neg \langle \text{non-majority} \rangle \omega_1) \land (\neg \langle \text{non-majority} \rangle \omega_2)$
    \item Guaranteed outcome: $\langle \text{any} \rangle (\omega_1 \lor \omega_2)$
    \item Incompatibility: $\langle \text{any} \rangle \neg (\omega_1 \land \omega_2)$
  \end{itemize}
\end{frame}

\note{
  + This is the crown jewel application: voting systems.
  + In QCL, you can specify a complete majority voting system in 4 short formulas.
  + The predicates capture the essence: ``any group of size > n/2 can determine the outcome.''
  + Compare this to CL, where each formula would need to list all majority coalitions explicitly.
  + For $n = 100$, the CL version would be astronomically long.
  + Timing: 26:00 to 28:00.
}

\begin{frame}{Why This Matters}
  \begin{table}
    \centering
    \begin{tabular}{lcc}
      \toprule
      \textbf{Property} & \textbf{CL} & \textbf{QCL} \\
      \midrule
      Expressive power & Full & Same as CL \\
      Formula length & Exponential & Polynomial \\
      Satisfiability & PSPACE-complete & PSPACE-complete \\
      Model checking & P / PSPACE & P / PSPACE \\
      \bottomrule
    \end{tabular}
  \end{table}
\end{frame}

\note{
  + Summary table comparing CL and QCL.
  + Same expressive power, same complexity, but QCL is exponentially more succinct.
  + This makes QCL the practical choice for specifying multi-agent systems.
  + The theoretical foundation (effectivity functions, game frames) remains the same.
  + Timing: 28:00 to 29:00.
}

% ============================================================
% SLIDE 21 — Conclusions
% ============================================================
\section{Conclusions and Reflections}

\begin{frame}{The Unifying Vision}
  \begin{itemize}
    \item Coalition Logic provides a \textbf{unifying game-theoretic view} of modal logic.
    \item Normal and non-normal modal logics emerge as 1- and 2-player versions of CL.
    \item Every formula is anchored to a concrete game frame.
  \end{itemize}
\end{frame}

\note{
  + Conclude by emphasizing the unifying power of CL.
  + Traditional modal logic is just the 1-player case of Coalition Logic.
  + QCL bridges theory and practice: you can now specify and verify real systems.
  + ``Social Software'' is Pauly's vision: using logic to design and verify social institutions.
  + Thank the audience. Open for questions (10 minutes).
  + Timing: 29:00 to 30:00.
}

\begin{frame}{From Theory to Practice}
  \begin{itemize}
    \item QCL makes complex social mechanisms (voting, cooperation) \textbf{specifiable and verifiable}.
    \item Marc Pauly's ``Social Software'': formally verifying the robustness of institutions.
    \item Future work: integrating preferences and knowledge for deeper understanding of coalition formation.
  \end{itemize}
\end{frame}

\begin{frame}
  \centering
  \Huge \textit{Questions?}
\end{frame}

\note{
  + Thank the audience for their attention.
  + Remind them of the 10-minute Q\&A session.
  + Be prepared to explain:
    - Why superadditivity makes CL PSPACE-complete.
    - The difference between $[P]\varphi$ and $\neg \langle P \rangle \neg \varphi$.
    - How to read speaker notes with pdfpc.
}

\end{document}
